\documentclass[11pt]{article}

\usepackage[margin=1.05in]{geometry}
\usepackage[T1]{fontenc}
\usepackage[utf8]{inputenc}
\usepackage{lmodern}
\usepackage{amsmath,amssymb,amsthm,mathtools}
\usepackage{hyperref}
\usepackage{enumitem}
\usepackage{xcolor}
\usepackage{booktabs}

\hypersetup{
    colorlinks=true,
    linkcolor=blue!50!black,
    citecolor=blue!50!black,
    urlcolor=blue!50!black
}

\newtheorem{theorem}{Theorem}[section]
\newtheorem{lemma}[theorem]{Lemma}
\newtheorem{proposition}[theorem]{Proposition}
\newtheorem{corollary}[theorem]{Corollary}
\theoremstyle{definition}
\newtheorem{definition}[theorem]{Definition}
\newtheorem{remark}[theorem]{Remark}

\newcommand{\cp}{\operatorname{cp}}
\newcommand{\CoreCap}{\operatorname{Cap}}
\newcommand{\calT}{\mathcal T}
\newcommand{\calR}{\mathcal R}
\newcommand{\calA}{\mathcal A}
\newcommand{\calL}{\mathcal L}
\newcommand{\rel}{\mathrm{rel}}
\newcommand{\own}{\operatorname{own}}

\title{A Fractional Triangle Theorem for Chordal Graphs}
\author{
Juan Pablo Traverso Gianini\\
Independent researcher, Santiago, Chile\\
\texttt{jtraverso@gmail.com}\\
\href{https://orcid.org/0009-0003-6068-4096}{ORCID: 0009-0003-6068-4096}
}
\date{Preprint v0.3 --- July 1, 2026}

\begin{document}

\maketitle

\begin{center}
{\small Paper C in the series on clique partitions in split and chordal graphs. Preprint; imports the finite boundary-export \(L_4\) certificate package from Paper B.}
\end{center}

\noindent\textbf{Keywords.}
Clique partitions; chordal graphs; clique trees; triangle packings; fractional packings; edge clique partitions; Erd\H{o}s problems.

\medskip

\noindent\textbf{MSC 2020.}
05C70, 05C35, 05C65, 05C69.

\begin{abstract}
We prove the fractional triangle-packing theorem needed for the asymptotic chordal clique partition bound, conditional on the local toolkit and finite exact \(L_4\) certificate package of Paper B. The proof uses a nearest-heavy laminar decomposition of a clique tree into anchored one-boundary branches and two-boundary corridors. Terminal regions are classified using the Boundary-Export Simplicial \(L_4\) Alternative from Paper B. Boundary interactions are converted into debits, allocated over retained cores, and realized as fractional triangles by a uniform allocation lemma. Version v0.3 adds an explicit edge-wise retained-core ownership ledger, restricts uniform debit realization to subcores of size at least two, imports the corrected weighted Heavy-Bag lemma, and separates feasibility ownership from quantitative telescoping.
\end{abstract}

\tableofcontents

\section{Introduction}

The edge clique partition number of a graph \(G\), denoted \(\cp(G)\), is the minimum number of cliques whose edge sets partition \(E(G)\). A triangle packing immediately gives a clique partition upper bound: starting from singleton-edge cliques, every edge-disjoint triangle can be replaced by one triangle clique, saving two cliques. Thus
\[
    \cp(G)\le |E(G)|-2\nu_3(G).
\]
For dense asymptotic questions it is natural to pass through fractional triangle packings. This paper proves the fractional chordal theorem needed for the chordal clique partition problem.

\subsection*{Role in the series}

The papers are organized as follows.

\begin{enumerate}[label=(\Alph*)]
    \item Paper A proves the sharp split graph theorem and supplies the lower-bound construction \(K_p\vee\overline{K}_{2p}\).
    \item Paper B proves the local and structural toolkit for chordal graphs: heavy bags, local reductions, \(L_4\) certificates, promotion lemmas, and the Core-Fan Allocation Theorem.
    \item Paper C, the present paper, proves the global fractional triangle theorem for chordal graphs using Paper B.
    \item Paper D applies fractional-to-integral triangle-packing rounding and derives the final chordal clique partition theorem.
\end{enumerate}

This paper does not prove the split lower bound and does not perform the final integral rounding.

\subsection*{Main theorem}

\begin{theorem}[Fractional Chordal Theorem]\label{thm:fractional-main}
Assuming the complete Paper B local toolkit in the form stated in Section~\ref{sec:imported-toolkit}, every chordal graph \(G\) on \(n\) vertices satisfies
\[
    |E(G)|-2\nu_3^*(G)
    \le
    \left(\frac16+o(1)\right)n^2.
\]
\end{theorem}

At the level of this preprint, the theorem is proved as an implication from the Paper B toolkit. The main new ingredients in this paper are the anchored recursive decomposition, the uniform realization of boundary debits, the retained-core ownership ledger, and the global amortization.

\section{Fractional triangle packings}

Let \(K_3(G)\) denote the set of triangles in \(G\). A fractional triangle packing is a choice of weights
\[
    \psi_T\ge0
    \qquad
    (T\in K_3(G))
\]
such that
\[
    \sum_{T\ni e}\psi_T\le1
    \qquad
    (e\in E(G)).
\]
Define
\[
    \nu_3^*(G)=
    \max\sum_{T\in K_3(G)}\psi_T.
\]
It is convenient to define the fractional triangle residual
\[
    \tau_3(G)=|E(G)|-2\nu_3^*(G).
\]
Theorem~\ref{thm:fractional-main} says
\[
    \tau_3(G)\le \left(\frac16+o(1)\right)n^2
\]
for chordal graphs.

\section{Chordal graphs and clique trees}

We use the subtree representation of chordal graphs. Let \(\calT\) be a clique tree. Each vertex \(v\in V(G)\) corresponds to a connected subtree
\[
    \calT_v\subseteq\calT.
\]
For a clique-tree node \(t\), set
\[
    C(t)=\{v:t\in\calT_v\}.
\]
Then
\[
    uv\in E(G)
    \quad\Longleftrightarrow\quad
    \calT_u\cap\calT_v\neq\varnothing.
\]
The bag coverage is
\[
    M(t)=|C(t)|.
\]

\section{Imported toolkit from Paper B}\label{sec:imported-toolkit}

We import the following statements from Paper B.

\begin{lemma}[Heavy-Bag Lemma, Paper B]\label{lem:import-heavy}
Let \(M_* = \max_t |C(t)|\). Then
\[
    |E(G)|\le M_* n.
\]
More generally, for nonnegative weights \(x_v\), define
\[
    M_x^*
    =
    \max_t\sum_{v:t\in\calT_v}x_v.
\]
Then
\[
    \sum_{uv\in E(G)}x_ux_v
    \le
    M_x^*\sum_vx_v.
\]
\end{lemma}

\begin{theorem}[Local Terminal Toolkit, Paper B]\label{thm:import-toolkit}
Every bounded local residue arising in the chordal recursion is good-leaf reducible, ear-budget reducible by bounded-multiplicity simplicial absorption, sealed or boundary-export \(L_4\)-certified with boundary rank at most two, promoted by a long-core or overload rule, assigned to parent-core Core-Fan/Farkas allocation, or charged to the lower-order residual ledger.
\end{theorem}

\begin{theorem}[Core-Fan Allocation Theorem, Paper B]\label{thm:import-core-fan}
For every family \(\calA\) of debit-active boundary groups, with \(s_\alpha=|S_\alpha|\ge2\), and debits
\[
    T_\alpha
    =
    \frac{u_\alpha(4s_\alpha-u_\alpha-2r_\alpha)_+}{12},
\]
one has
\[
    \sum_{\alpha\in\calA}T_\alpha
    \le
    \CoreCap
    \left(
    \bigcup_{\alpha\in\calA}E(S_\alpha)
    \right).
\]
\end{theorem}

\begin{theorem}[Boundary-Export \(L_4\) Certificate Theorem, Paper B]\label{thm:import-l4}
Every bounded span-four \(L_4\) residue with boundary rank at most two admits a rational type-level fractional triangle certificate. Internal and mixed triangles are locally owned. Two-boundary-one-local triangles export their boundary-boundary load as core debits. No pure-boundary triangle is consumed locally. The sealed certificate is the rank-zero case.
\end{theorem}

\begin{lemma}[Type Blow-Up Realization, Paper B]\label{lem:import-blowup}
A rational type-level fractional triangle certificate over a fixed type system lifts to an actual fractional triangle packing with \(O(n)\) loss.
\end{lemma}

\begin{remark}[Interface status]
The present preprint uses the Paper B local toolkit and finite boundary-export \(L_4\) certificate family as imported components. The certificate data, verifier, and verifier output must accompany Paper B for independent reproducibility.
\end{remark}

\section{Anchored regions}

The recursive proof works with anchored regions in the clique tree.

\begin{definition}[One-boundary branch]
A one-boundary region is a pair
\[
    \calR=(R;S),
\]
where \(R\subseteq\calT\) is a connected subtree and \(S\) is a retained boundary clique or subcore. Its private set is
\[
    P(\calR)
    =
    \{v:\calT_v\cap R\neq\varnothing\}\setminus S.
\]
Its relative edge set is
\[
    E^{\rel}(\calR)
    =
    E(G[P(\calR)\cup S])\setminus E(G[S]).
\]
\end{definition}

\begin{definition}[Two-boundary corridor]
A two-boundary region is
\[
    \calR=(R;S_L,S_R),
\]
with two retained boundary cores. Its private set is
\[
    P(\calR)
    =
    \{v:\calT_v\cap R\neq\varnothing\}
    \setminus(S_L\cup S_R),
\]
and its relative edge set is
\[
    E^{\rel}(\calR)
    =
    E(G[P(\calR)\cup S_L\cup S_R])
    \setminus
    \left(E(G[S_L])\cup E(G[S_R])\right).
\]
\end{definition}

Let
\[
    m(\calR)=|P(\calR)|.
\]
For a clique-tree node \(t\) internal to \(R\), define private coverage
\[
    M_{\calR}(t)
    =
    |\{v\in P(\calR):t\in\calT_v\}|.
\]

Fix \(\eta>0\). A region is \(\eta\)-diffuse if
\[
    M_{\calR}^*
    :=
    \max_{t\in R^\circ}M_{\calR}(t)
    \le
    \eta m(\calR).
\]

\section{Nearest-heavy laminar recursion}

If \(\calR\) is not diffuse, choose a nearest heavy bag \(b\in R^\circ\) satisfying
\[
    M_{\calR}(b)>\eta m(\calR).
\]
Retain
\[
    C_b=C(b)
\]
as a new core and split the tree region \(R\) along the node \(b\). Components on old boundary sides become two-boundary corridors; off-path components become one-boundary branches attached to \(C_b\).

\begin{theorem}[Nearest-Heavy Laminar Decomposition]\label{thm:nearest-heavy}
For fixed \(\eta>0\), the nearest-heavy recursion produces a finite laminar family of anchored terminal regions and retained cores such that:
\begin{enumerate}[label=(\roman*)]
    \item terminal private sets are pairwise disjoint;
    \item retained core edges are excluded from child-private accounts;
    \item diffuse terminal errors sum to \(O(\eta n^2)\);
    \item potential releases telescope under \(F(x)=x^2/6\).
\end{enumerate}
\end{theorem}

\begin{proof}
Removing a heavy node \(b\) from the tree region \(R\) separates \(R\) into connected components. A component can be adjacent to the new core \(C_b\) and to at most one old boundary side, since the path between two old boundary sides passes through \(b\). Hence every child is again a one-boundary or two-boundary region.

The number of internal clique-tree nodes strictly decreases along every child, because \(b\) becomes retained core boundary and is no longer internal. Thus the recursion terminates.

For disjointness, suppose a vertex \(v\) were private in two different child components. Since \(\calT_v\) is connected and intersects both components of \(R-b\), it must contain \(b\), hence \(v\in C_b\). But vertices of \(C_b\) are retained as boundary/core vertices in the children, not private vertices. Contradiction.

The same laminarity gives global disjointness of terminal private masses:
\[
    \sum_{\calR\in\calL}m(\calR)\le n.
\]
Diffuse terminal regions contribute \(O(\eta m(\calR)^2)\), hence in total
\[
    \sum_{\calR\in\calL}O(\eta m(\calR)^2)
    \le
    O(\eta n^2).
\]
Finally, because private masses separated at different recursion nodes are disjoint and retained cores are not deleted as child-private mass, the potential drops telescope under \(F(x)=x^2/6\).
\end{proof}

\section{Retained-core ownership ledger}\label{sec:ownership-ledger}

The recursion may create retained cores that overlap as vertex sets. The proof does not require retained cores to be disjoint as sets of vertices. It requires that every edge capacity be owned by a unique active account.

\begin{lemma}[Retained-core ownership ledger]\label{lem:retained-core-ledger}
Run the nearest-heavy recursion with the retained-core convention. Then there is a single-valued ownership map
\[
    \omega:E(G)\to
    \{\mathrm{internal},\mathrm{boundary},\mathrm{core},\mathrm{unused}\}
\]
with the following properties.
\begin{enumerate}[label=(\roman*)]
    \item If both endpoints of an edge are private in some recursive region, then the edge is owned by the deepest region in which both endpoints are private.
    \item If one endpoint is private in a region and the other lies in its retained boundary/core, then the edge is owned by the boundary-debit account of that region, unless it is already internally owned by a deeper region.
    \item If both endpoints lie in one or more retained cores, then the edge is owned by the deepest retained core containing both endpoints.
    \item Retained-core edges are excluded from all descendant relative edge sets.
    \item Every boundary debit supported on a subcore \(S_\alpha\) is charged only to the retained core owning the edges of \(E(S_\alpha)\).
    \item No edge is used by two different local or boundary triangle-packing accounts.
\end{enumerate}
\end{lemma}

\begin{proof}
The recursion is laminar on clique-tree regions. At each non-diffuse step, a heavy bag \(b\) is retained as a core and the remaining components of the clique-tree region become child regions. These child regions are pairwise disjoint as tree regions, except for their retained boundary/core vertices.

Suppose a vertex \(v\) were private in two different child regions after splitting at \(b\). Since the clique-tree support \(\calT_v\) is connected and meets two distinct components of \(R-b\), it must contain \(b\). Hence \(v\in C_b\), so \(v\) is retained as a boundary/core vertex, not private in either child. Thus child private sets are disjoint. By induction, terminal private sets are disjoint globally.

Now fix an edge \(xy\in E(G)\). Consider the set of recursive regions in which both \(x\) and \(y\) are private. Since the recursion tree is laminar, this set is linearly ordered by inclusion. If it is nonempty, choose the deepest such region \(R\). Then \(xy\) is owned by the internal account of \(R\). No ancestor may also use it, and no descendant contains both endpoints privately by maximality.

If no region contains both endpoints privately, but \(x\) is private in a region \(R\) and \(y\in B(R)\), then the edge \(xy\) is a private-boundary edge. It is assigned to the boundary-debit account of \(R\). It cannot be used by a child internal packing because \(y\) is retained boundary/core, not private.

Finally, if both endpoints lie in retained cores, choose the deepest retained core containing both endpoints. That core owns \(xy\). By the retained-core convention, core edges are removed from all descendant relative edge sets, so descendants cannot use \(xy\) internally. Boundary debits using \(xy\) report edge-wise to the owner of \(xy\); no assumption is made that all edges of a boundary subcore have the same retained-core owner.

Therefore every edge has at most one active owner. Since every local or boundary triangle-packing source is required to use only edges owned by its account, no edge is loaded by two different accounts.
\end{proof}

\begin{remark}[Multi-core bookkeeping]
This lemma is the explicit version of the multi-core bookkeeping used informally in earlier versions. Nested retained cores may overlap on vertices, but edge capacity is charged to the deepest retained core containing both endpoints. Descendant regions never reuse retained-core edges as private capacity.
\end{remark}


\section{Terminal classification}

Paper B supplies the Boundary-Export Simplicial \(L_4\) Alternative.

\begin{theorem}[Terminal Classification]
\label{thm:terminal-classification}
Every terminal region produced by Theorem~\ref{thm:nearest-heavy} is one of:
\begin{enumerate}[label=(\roman*)]
    \item an \(\eta\)-diffuse one-boundary branch;
    \item an \(\eta\)-diffuse two-boundary corridor;
    \item an atomic core-fan debit call;
    \item a sealed or boundary-export \(L_4\)-certified window with boundary rank at most two;
    \item a good-leaf or bounded-multiplicity ear-budget reducible window;
    \item a common-core/sunflower or long-core/five-window promotion event;
    \item parent-core allocation or lower-order residue.
\end{enumerate}
\end{theorem}

\begin{proof}
If a terminal region is diffuse, it is of type (i) or (ii). Suppose it is not diffuse. If it had an internal heavy node, the nearest-heavy recursion would split it further. Thus a non-diffuse terminal region has no unresolved heavy split in the sense of the recursion. Apply the Paper B local terminal toolkit. Good endpoints reduce by good-leaf moves; bounded-multiplicity ears are absorbed locally; multiplicity-five overloads promote; boundary rank at most two is discharged by boundary-export \(L_4\) with exported core debits; boundary rank at least three is parent-core involvement and is promoted or globally allocated. These alternatives are exactly (iii)--(vii).
\end{proof}

\begin{remark}[No unrestricted window sweep]
The classification does not use an unrestricted overlapping \(L_4\) window-sweep lemma. The \(L_4\) certificate is invoked only after the budgeted sweep and promotion alternatives have reduced the residue to a sealed or boundary-export instance of boundary rank at most two.
\end{remark}

\section{Boundary debits}

For a boundary trace group \(\alpha\), let \(u_\alpha\) be its private mass, \(S_\alpha\) its retained subcore, \(s_\alpha=|S_\alpha|\), and \(r_\alpha\) its exterior mass.

\begin{definition}[Active boundary debit]
Define
\[
T_\alpha
=
\begin{cases}
\dfrac{u_\alpha(4s_\alpha-u_\alpha-2r_\alpha)_+}{12},
& s_\alpha\ge2,\\[1.2ex]
0,
& s_\alpha\le1.
\end{cases}
\]
A group with \(s_\alpha\ge2\) is debit-active. Groups with \(s_\alpha\le1\) do not enter the quadratic core-capacity ledger.
\end{definition}

By the imported Core-Fan Allocation Theorem, for every family \(\calA\) of debit-active groups,
\[
    \sum_{\alpha\in\calA}T_\alpha
    \le
    \CoreCap
    \left(
    \bigcup_{\alpha\in\calA}E(S_\alpha)
    \right).
\]
Thus boundary debits satisfy all Hall-type capacity inequalities over retained cores.

\begin{remark}[One-vertex and zero-vertex subcores]
If \(s_\alpha\le1\), then \(S_\alpha\) contains no core edge. Therefore there is no edge \(ij\subseteq S_\alpha\) on which to place a triangle \(uij\). Such groups are left unused or charged to the lower-order/diffuse residual ledger according to the terminal classification. This only weakens the fractional packing and prevents an undefined division by \(\binom{s_\alpha}{2}\).
\end{remark}

\section{Uniform boundary-debit realization}

Paper B proves that debit-active boundary debits fit in core capacity. Paper C must realize them as actual fractional triangles.

For every debit-active group \(\alpha\), define
\[
    z_{\alpha,ij}
    =
    \frac{T_\alpha}{\binom{s_\alpha}{2}},
    \qquad
    ij\subseteq S_\alpha.
\]
No such quantity is defined for \(s_\alpha\le1\).

\begin{lemma}[Uniform Allocation Calculus]\label{lem:uniform-calculus}
For \(v_i\ge0\), \(V=\sum_i v_i<1\), and \(0<q_i\le1-V\),
\[
\sum_i
\frac{
v_i(6q_i+v_i-2)_+
}{6q_i^2}
\le
\frac23.
\]
\end{lemma}

\begin{proof}[Proof sketch]
The summand is zero unless \(v_i>2-6q_i\). On each active interval it is convex in the relevant normalized variable. Under the constraints \(v_i\ge0\), \(\sum_i v_i=V\), and \(q_i\le1-V\), the maximum occurs at an extreme distribution. The one-variable extreme case reduces to
\[
    \frac{v(6q+v-2)}{6q^2}
    \le
    \frac23,
\]
with equality only in the limiting balanced fan case. Summing over the extremal active groups gives the displayed bound.
\end{proof}

\begin{lemma}[Uniform Boundary-Debit Realization]\label{lem:uniform-realization}
The allocation \(z_{\alpha,ij}\), restricted to debit-active groups \(s_\alpha\ge2\), realizes all active boundary debits by fractional triangles while respecting private-boundary and core-edge capacities.
\end{lemma}

\begin{proof}
For a debit-active group \(\alpha\), distribute \(z_{\alpha,ij}\) equally over core edges \(ij\subseteq S_\alpha\). Then
\[
    \sum_{ij\subseteq S_\alpha}z_{\alpha,ij}=T_\alpha.
\]
Each unit of \(z_{\alpha,ij}\) is realized by triangles whose private vertex lies in group \(\alpha\) and whose two core vertices are \(i,j\). The private-boundary load incident to a fixed private vertex is bounded by the uniform fan calculation, and for a core edge \(e\), the total normalized load is bounded by Lemma~\ref{lem:uniform-calculus}. The Core-Fan Hall/Farkas inequalities ensure that the aggregate requested core-edge load lies within retained-core capacity.

Groups with \(s_\alpha\le1\) have \(T_\alpha=0\) by definition and are not realized by core-edge triangles. Hence the allocation is well-defined and all active debits are realized.
\end{proof}


\section{Boundary-export \(L_4\) windows and type blow-up}

Paper B supplies the Boundary-Export \(L_4\) Certificate Theorem and Type Blow-Up Realization Lemma. Thus sealed or boundary-export \(L_4\) terminal windows contribute only \(O(|X|)\) local realization loss, and over disjoint private interiors this sums to
\[
    O_\eta(n).
\]

In boundary-export mode, triangles with no boundary vertex and triangles with one boundary vertex are locally owned. Triangles with two boundary vertices and one local vertex export the boundary-boundary edge load as a debit to the retained parent core; those debits enter the Core-Fan/Farkas ledger. Pure-boundary triangles are not consumed locally. The sealed mode is the special case of boundary rank zero.

\section{Edge ownership patching}

\begin{theorem}[Edge Ownership Patching Theorem]\label{thm:ownership-patching}
Suppose every local, sealed, internal, and boundary fractional triangle assignment obeys the ownership map of Lemma~\ref{lem:retained-core-ledger}. Then the sum of all local assignments is a feasible global fractional triangle packing of \(G\).
\end{theorem}

\begin{proof}
Let \(e\in E(G)\). By Lemma~\ref{lem:retained-core-ledger}, \(e\) has at most one active ownership account. All triangle weights produced by an account are required to use only edges owned by that account and to respect capacity one on those edges. Since no other account can use \(e\), the total global load on \(e\) is at most one. This holds for every edge, so the sum of all local and boundary assignments is a feasible fractional triangle packing.
\end{proof}

\section{Lower-order ledger}

\begin{lemma}[Lower-Order Ledger]\label{lem:lower-order}
All finite, sealed, and type-blow-up losses are \(O_\eta(n)\). Diffuse terminal losses are \(O(\eta n^2)\). Boundary groups with \(s_\alpha\le1\) contribute no active quadratic debit and are assigned to the unused or residual ledger. Therefore, after summing over the laminar terminal family, the total unaccounted loss is
\[
    O(\eta n^2)+O_\eta(n).
\]
\end{lemma}

\begin{proof}
Terminal private sets are pairwise disjoint by Theorem~\ref{thm:nearest-heavy}. Every fixed-size certificate or type-blow-up loss is linear in the private mass of its terminal region, so the sum of all such losses is \(O_\eta(n)\). Diffuse terminal regions have maximum private bag coverage at most \(\eta m(\calR)\); by the Heavy-Bag Lemma applied inside the private region, their remaining internal interaction is \(O(\eta m(\calR)^2)\). Summing over disjoint terminal private masses gives \(O(\eta n^2)\).

Groups with \(s_\alpha\le1\) do not request core-edge capacity and hence cannot create double-counted core load. Their edges are either covered by an internal local packing or left unused. Leaving them unused can only increase the residual by the amount already assigned to the terminal residual ledger. Thus they do not affect the quadratic core-fan capacity calculation.
\end{proof}

\section{Proof of the fractional theorem}

\begin{proof}[Proof of Theorem~\ref{thm:fractional-main}]
Fix \(\eta>0\). Apply the nearest-heavy laminar recursion to the clique tree. Terminal regions are classified by Theorem~\ref{thm:terminal-classification}. Diffuse regions contribute \(O(\eta n^2)\) residual. Good-leaf and ear-budget regions are locally reducible. Sealed or boundary-export \(L_4\) windows are packed using the imported certificate and type blow-up realization, with total loss \(O_\eta(n)\). Exported \(L_4\) debits and ordinary debit-active boundary groups report debits to their retained cores; the Core-Fan Allocation Theorem gives capacity feasibility, and Lemma~\ref{lem:uniform-realization} realizes the debits as fractional triangles.

By Lemma~\ref{lem:retained-core-ledger} and Theorem~\ref{thm:ownership-patching}, all local and boundary triangle weights patch into one feasible global fractional triangle packing. The potential releases telescope under \(F(x)=x^2/6\), and the lower-order ledger contributes
\[
    O(\eta n^2)+O_\eta(n).
\]
Therefore
\[
    |E(G)|-2\nu_3^*(G)
    \le
    \frac{n^2}{6}+O(\eta n^2)+O_\eta(n).
\]
Letting \(n\to\infty\) first and then \(\eta\to0\) gives
\[
    |E(G)|-2\nu_3^*(G)
    \le
    \left(\frac16+o(1)\right)n^2.
\]
\end{proof}

\section{Relation to integral clique partitions}

The final integral theorem is not proved in this paper. Paper D uses the fixed-family rounding theorem of Haxell--R\"odl and Yuster:
\[
    \nu_3(G)
    \ge
    \nu_3^*(G)-o(n^2)
\]
for the fixed family \(\{K_3\}\). Hence
\[
    \cp(G)
    \le
    |E(G)|-2\nu_3(G)
    \le
    |E(G)|-2\nu_3^*(G)+o(n^2).
\]
Combining this with Theorem~\ref{thm:fractional-main} yields the chordal clique partition upper bound.

\section{Discussion and further directions}

Theorem~\ref{thm:fractional-main} is the conceptual center of the series. It separates the chordal clique partition problem into two layers: a fractional triangle-packing statement proved by clique-tree recursion, and a final integral rounding step handled in Paper D.

Version v0.3 stabilizes three reviewer-sensitive interfaces:
\begin{enumerate}[label=(\roman*)]
    \item the weighted Heavy-Bag import now uses weighted bag maximum \(M_x^*\);
    \item uniform boundary-debit realization is restricted to subcores \(s_\alpha\ge2\);
    \item the retained-core ownership ledger is promoted to a main lemma, rather than hidden as lower-order bookkeeping.
\end{enumerate}

The remaining reproducibility requirement is the finite Paper B boundary-export \(L_4\) certificate package. The global ownership assembly, telescoping, and debit realization are explicit in this preprint.

\section{Proof-status audit}

\begin{center}
\begin{tabular}{lll}
\toprule
Item & Location & Status in v0.3\\
\midrule
Nearest-heavy recursion & Paper C & proof complete\\
Retained-core ownership ledger & Paper C & explicit proof added\\
Terminal classification & Paper C + Paper B & closed modulo finite \(L_4\) data\\
Core-Fan capacity & Paper B & imported for debit-active groups\\
Uniform allocation & Paper C & guarded by \(s_\alpha\ge2\)\\
Boundary-export \(L_4\) certificate & Paper B & imported finite exact package\\
Type blow-up & Paper B & imported\\
Lower-order ledger & Paper C & proof complete\\
Fractional theorem & Paper C & conditional assembly\\
Integral rounding & Paper D & not included here\\
\bottomrule
\end{tabular}
\end{center}

\section*{Acknowledgements and use of AI-assisted tools}

The author is deeply grateful to his wife Mar\'ia Paz and to his children Lucas, Juan Crist\'obal, Francisca, Raimundo, and Benjam\'in for their love, patience, and support during the development of this work.

The author used AI-assisted tools, including Gemini 3.5 Pro and ChatGPT, during the development of this work for exploratory search, testing possible approaches, looking for contradictions, organizing proof strategies, improving exposition, and preparing drafts. All mathematical claims, proofs, citations, and final editorial decisions are the sole responsibility of the author. No AI system is listed as an author.

\begin{thebibliography}{9}

\bibitem{ChenErdosOrdman}
G.-T. Chen, P. Erd\H{o}s, and E. T. Ordman,
\emph{Clique partitions of split graphs},
in \emph{Combinatorics, graph theory, algorithms and applications} (Beijing, 1993), 21--30, 1994.

\bibitem{ErdosOrdmanZalcstein}
P. Erd\H{o}s, E. T. Ordman, and Y. Zalcstein,
\emph{Clique partitions of chordal graphs},
Combinatorics, Probability and Computing (1993), 409--415.

\bibitem{EGP66}
P. Erd\H{o}s, A. W. Goodman, and L. P\'osa,
\emph{The representation of a graph by set intersections},
Canadian Journal of Mathematics 18 (1966), 106--112.

\bibitem{Gavril}
F. Gavril,
\emph{The intersection graphs of subtrees in trees are exactly the chordal graphs},
Journal of Combinatorial Theory, Series B 16 (1974), 47--56.

\bibitem{HaxellRodl}
P. E. Haxell and V. R\"odl,
\emph{Integer and fractional packings in dense graphs},
Combinatorica 21 (2001), 13--38.

\bibitem{Yuster}
R. Yuster,
\emph{Integer and fractional packing of families of graphs},
Combinatorica 28 (2008), 245--251.

\end{thebibliography}

\end{document}
